% Arquitetura de Inmon: de cima para baixo (cap. 2).
\begin{tikzpicture}[diagrama, x=1mm, y=1mm]
  % sistemas de origem
  \foreach \y/\nome/\id in {13.5/ERP/f1, 4.5/CRM/f2, -4.5/Planilhas/f3, -13.5/{Outros}/f4}
    \node[caixa, minimum width=15mm, minimum height=6mm] (\id) at (8,\y) {\nome};
  \node[moldura, fit=(f1)(f4), label={[rotulo]above:Sistemas de origem}] {};

  % armazem corporativo
  \node[bancodestaque, minimum width=26mm, minimum height=26mm] (edw) at (46,0)
    {Data warehouse\\corporativo\\[1mm]{\scriptsize normalizado (3FN)}\\{\scriptsize e integrado}};
  \foreach \id in {f1,f2,f3,f4} \draw[seta] (\id.east) -- (edw.west);
  \node[rotulo] at (25.5,18) {ETL};

  % data marts dependentes
  \foreach \y/\nome/\id in {13/Vendas/m1, 0/Finanças/m2, -13/Logística/m3}
    \node[banco, minimum width=17mm, minimum height=9mm] (\id) at (79,\y) {\nome};
  \node[moldura, fit=(m1)(m3), label={[rotulo]above:Data marts dependentes}] {};
  \foreach \id in {m1,m2,m3} \draw[seta] (edw.east) -- (\id.west);

  % consumo
  \foreach \y/\nome/\id in {13/Relatórios/u1, 0/Dashboards/u2, -13/Análises/u3}
    \node[caixa, minimum width=16mm, minimum height=6mm] (\id) at (108,\y) {\nome};
  \foreach \i in {1,2,3} \draw[seta] (m\i.east) -- (u\i.west);

  \node[rotulo, text width=110mm] at (58,-27)
    {Ordem de construção: primeiro o armazém corporativo, depois os data marts de cada área.};
\end{tikzpicture}
